\documentclass[11pt]{article}
\usepackage[margin=1in]{geometry}
\usepackage{enumitem}
% =============================================================================
% Preambles/header.tex — shared LaTeX/Beamer preamble for this project
%
% Use from a Beamer lecture by prepending:
%     \input{header}
% and compiling with TEXINPUTS=../Preambles:$TEXINPUTS (the /compile-latex
% skill does this automatically).
%
% PALETTE CONTRACT. Color names defined here MUST match the names in
% Quarto/theme-template.scss so Beamer and Quarto renderings use the same
% palette. The scripts/check-palette-sync.sh script enforces this.
%
% To customize: edit the HEX values below AND the matching $variable in
% Quarto/theme-template.scss, then run scripts/check-palette-sync.sh.
% =============================================================================

% -----------------------------------------------------------------------------
% Engine detection + encoding
%
% Our /compile-latex skill uses XeLaTeX, where inputenc is unnecessary and
% can produce encoding warnings. Only load inputenc under pdfTeX.
% -----------------------------------------------------------------------------
\usepackage{iftex}
\ifPDFTeX
  \usepackage[utf8]{inputenc}
\fi

\usepackage{amsmath, amssymb, amsthm}
\usepackage{booktabs}
\usepackage{graphicx}

% -----------------------------------------------------------------------------
% PALETTE — mirrors Quarto/theme-template.scss variable names.
% Load xcolor and define colors BEFORE anything that references them
% (notably hyperref, hypersetup, and the Beamer color assignments).
% -----------------------------------------------------------------------------
\usepackage{xcolor}

\definecolor{primary-blue}{HTML}{012169}      % headings, accents
\definecolor{primary-gold}{HTML}{B9975B}      % emphasis, borders
\definecolor{highlight-yellow}{HTML}{F2A900}  % markers, alerts
\definecolor{light-bg}{HTML}{E8EDF5}          % title / section backgrounds
\definecolor{jet}{HTML}{1A1A1A}               % body text

% Semantic colors (used in diagrams and annotations)
\definecolor{positive}{HTML}{15803D}          % good / observed / identified
\definecolor{negative}{HTML}{B91C1C}          % bad / problematic / bias
\definecolor{neutral}{HTML}{525252}           % reference / context

% Highlight accents (match .hi-* classes in the SCSS)
\definecolor{hi-slate}{HTML}{314F4F}
\definecolor{hi-green}{HTML}{2E8B57}
\definecolor{hi-red}{HTML}{C41E3A}

% Semantic aliases so skill templates and snippets can write
% \textcolor{accent}{...} without hard-coding "primary-blue".
\colorlet{accent}{primary-blue}
\colorlet{accent2}{primary-gold}

% -----------------------------------------------------------------------------
% Hyperref — loaded AFTER xcolor + palette so link colors resolve.
% -----------------------------------------------------------------------------
\usepackage{hyperref}
\hypersetup{colorlinks=true, linkcolor=primary-blue, urlcolor=primary-blue,
            citecolor=primary-blue, pdfborder={0 0 0}}

% -----------------------------------------------------------------------------
% Course code — set once per deck (after \input{header}, before \begin{document})
% via \coursecode{CS401}. Shown in the footer on every slide so a deck's course
% is identifiable without opening the title page. Empty by default.
% -----------------------------------------------------------------------------
\makeatletter
\newcommand{\coursecode}[1]{\gdef\@coursecodetext{#1}}
\gdef\@coursecodetext{}
\makeatother

% -----------------------------------------------------------------------------
% Beamer theme assignments (only apply under Beamer).
%
% We detect Beamer via \@ifclassloaded{beamer}, which is the documented,
% non-internal way. This must be wrapped in \makeatletter/\makeatother
% because \@ifclassloaded contains an @-catcode character.
% -----------------------------------------------------------------------------
\makeatletter
\@ifclassloaded{beamer}{%
  \usefonttheme{professionalfonts}
  \setbeamercolor{structure}{fg=primary-blue}
  \setbeamercolor{frametitle}{fg=primary-blue}
  \setbeamercolor{title}{fg=primary-blue}
  \setbeamercolor{subtitle}{fg=primary-gold}
  \setbeamercolor{author}{fg=primary-blue}
  \setbeamercolor{institute}{fg=neutral}
  \setbeamercolor{date}{fg=primary-gold}
  \setbeamercolor{itemize item}{fg=primary-blue}
  \setbeamercolor{itemize subitem}{fg=primary-gold}
  \setbeamercolor{alerted text}{fg=primary-gold}
  \setbeamercolor{block title}{fg=white, bg=primary-blue}
  \setbeamercolor{block body}{bg=light-bg}
  % Semantic block variants: block=definition/concept (blue), exampleblock=
  % worked example (green/positive), alertblock=key takeaway or caution (gold).
  % Keep to <=2 colored boxes per slide across ALL three variants (INV-7).
  \setbeamercolor{block title example}{fg=white, bg=positive}
  \setbeamercolor{block body example}{bg=light-bg}
  \setbeamercolor{block title alerted}{fg=white, bg=primary-gold}
  \setbeamercolor{block body alerted}{bg=light-bg}

  % Minimal footer: course code (left) + page X / N (right), no navigation clutter
  \setbeamertemplate{navigation symbols}{}
  \setbeamertemplate{footline}{%
    \color{neutral}\footnotesize\hspace{0.7em}\@coursecodetext\hfill
    \insertframenumber/\inserttotalframenumber\hspace{0.7em}\vspace{0.3em}}
}{}
\makeatother

% -----------------------------------------------------------------------------
% TikZ setup — libraries the snippet gallery uses
% -----------------------------------------------------------------------------
\usepackage{tikz}
\usetikzlibrary{arrows.meta, positioning, calc, decorations.pathreplacing,
                fit, shapes.geometric, backgrounds}

% Shared TikZ styles — snippets and hand-written diagrams can pick these up
% via `\begin{tikzpicture}[every node/.style={...}]` or by naming specific
% styles (e.g., `\node[dag-node] (X) at (0,0) {X};`).
\tikzset{
  % Node styles for causal DAGs and similar
  dag-node/.style = {
    draw=primary-blue, fill=light-bg, rounded corners=2pt,
    minimum width=1.2cm, minimum height=0.9cm, align=center, font=\small
  },
  decision-node/.style = {
    draw=primary-gold, fill=white, diamond, aspect=1.8,
    minimum width=1.8cm, minimum height=1.2cm, align=center, font=\small,
    inner sep=1pt
  },
  % Edge styles — observed vs counterfactual
  observed-edge/.style = {-{Latex[length=2.5mm]}, thick, draw=jet},
  counterfactual-edge/.style = {-{Latex[length=2.5mm]}, thick, draw=neutral, dashed},
  confound-edge/.style = {-{Latex[length=2.5mm]}, thick, draw=negative, dashed},
  % Data-point styles for plots
  observed-dot/.style = {fill=primary-blue, draw=primary-blue, circle, inner sep=1.2pt},
  counterfactual-dot/.style = {fill=white, draw=neutral, circle, inner sep=1.2pt}
}

% -----------------------------------------------------------------------------
% Convenience macros
% -----------------------------------------------------------------------------
\newcommand{\muted}[1]{\textcolor{neutral}{#1}}
\newcommand{\key}[1]{\textcolor{primary-gold}{\textbf{#1}}}
\newcommand{\good}[1]{\textcolor{positive}{#1}}
\newcommand{\bad}[1]{\textcolor{negative}{#1}}

% Standout frame macro: full-bleed dark background for section transitions.
% Beamer-only (uses \begin{frame}/\setbeamercolor/\usebeamerfont) -- do not
% call from an article-class file (e.g. Notes/<CODE>/*-notes.tex). \muted,
% \key, \good, \bad, and \coursecode above are plain xcolor/gdef and are
% safe to use from either class; verified when Notes/article-class support
% was added.
% Expand: \transitionslide{Your transition title}
\newcommand{\transitionslide}[1]{%
  \begingroup
  \setbeamercolor{background canvas}{bg=primary-blue}
  \begin{frame}[plain]
    \centering
    \vfill
    {\usebeamerfont{title}\color{white}\Large #1\par}
    \vfill
  \end{frame}
  \endgroup
}

% Section divider frame (Beamer-only), an alternative to \transitionslide with a
% darker two-line style: near-black (jet) background, a small white label line
% above a larger gold title, and a thin gold rule along the bottom edge.
% Expand: \sectiondivider{Section 2}{Pointers}
\newcommand{\sectiondivider}[2]{%
  \begingroup
  \setbeamercolor{background canvas}{bg=jet}
  \begin{frame}[plain]
    \vspace*{3.0cm}
    {\color{white}\ttfamily\large #1\par}
    \vspace{0.5cm}
    {\color{primary-gold}\LARGE #2\par}
    \begin{tikzpicture}[remember picture, overlay]
      \draw[primary-gold, line width=2.5pt]
        ($(current page.south west)+(0,0.1)$) -- ($(current page.south east)+(0,0.1)$);
    \end{tikzpicture}
  \end{frame}
  \endgroup
}

% -----------------------------------------------------------------------------
% Channel watermark — "Concepts That Click" (YouTube).
%
% Article-class files (CompetitiveExam Q/A sets, Notes, Assignments) get a
% light diagonal draft watermark on every page plus a centered footer naming
% the channel. Beamer decks get a subtle TikZ background watermark instead
% (the footline already carries the course code). Centralized here so every
% MCQ PDF is branded without per-file edits.
% -----------------------------------------------------------------------------
\makeatletter
\@ifclassloaded{article}{%
  \usepackage{draftwatermark}
  \SetWatermarkText{Concepts That Click}
  \SetWatermarkScale{1}
  \SetWatermarkFontSize{2cm}
  \SetWatermarkLightness{0.86}
  \SetWatermarkAngle{45}
  \usepackage{fancyhdr}
  \pagestyle{fancy}
  \fancyhf{}
  \fancyfoot[C]{\footnotesize\textcolor{neutral}{Concepts That Click~|~YouTube: \href{https://www.youtube.com/@concepts.thatclick}{@concepts.thatclick}}}
  \renewcommand{\headrulewidth}{0pt}
  \renewcommand{\footrulewidth}{0pt}
  \fancypagestyle{plain}{%
    \fancyhf{}%
    \fancyfoot[C]{\footnotesize\textcolor{neutral}{Concepts That Click~|~YouTube: \href{https://www.youtube.com/@concepts.thatclick}{@concepts.thatclick}}}%
    \renewcommand{\headrulewidth}{0pt}%
    \renewcommand{\footrulewidth}{0pt}%
  }
}{}
\@ifclassloaded{beamer}{%
  \setbeamertemplate{background}{%
    \begin{tikzpicture}[remember picture, overlay]
      \node[rotate=45, opacity=0.055, align=center]
        at (current page.center)
        {\Huge\textcolor{neutral}{Concepts That Click}};
    \end{tikzpicture}%
  }
}{}
\makeatother 
\coursecode{JKSSB}
\renewcommand{\thesection}{18.\arabic{section}}
\newtheorem{example}{Example}[section]
\newenvironment{definitionbox}[1]{%
  \par\noindent\textbf{Definition (#1).}\ \itshape
}{\par\vspace{0.3em}}
\title{Software Categories --- Detailed Lecture Notes}
\author{JKSSB Computer Awareness}
\date{\today}

\begin{document}
\maketitle

\section{What is software?}
A computer system has two parts: hardware and software. \textbf{Hardware} is
the physical part that you can see and touch, such as the CPU, memory, and
input/output devices. \textbf{Software} is the set of programs, that is, the
instructions that tell the computer what to do. Programs are stored in memory
and executed by the CPU, and without software the hardware cannot perform any
useful task. Where hardware is the body of the computer, software is the set
of instructions that gives it purpose.

\begin{definitionbox}{Software}
A set of programs or instructions that directs the computer to perform a
task. Hardware is the physical equipment; software is the logical
instructions that run on it.
\end{definitionbox}

Software does not exist as a physical object; it exists as information. A
programmer writes the instructions in a programming language, and these
instructions are then translated, stored in memory, and executed by the CPU.
The same piece of hardware --- a desktop computer, for example --- can behave
completely differently depending on which software is loaded onto it. This is
why software is described as the logical part of a computer system: it
supplies the behaviour that the physical parts only make possible.

\begin{center}
\begin{tabular}{p{0.20\textwidth}p{0.32\textwidth}p{0.36\textwidth}}
\toprule
\textbf{Point} & \textbf{Hardware} & \textbf{Software} \\
\midrule
Nature & Physical; can be touched & Logical; a set of instructions \\
Made of & Electronic and mechanical components & Programs and data \\
Example & CPU, RAM, keyboard, monitor & Operating system, MS Word, a browser \\
Failure & Fails due to wear, damage, or age & Fails due to errors (bugs) or corruption \\
\bottomrule
\end{tabular}
\end{center}

Software is commonly grouped by \emph{what it is for} into six categories:
system software, application software, utility software, programming
software, embedded software, and middleware. These categories describe purpose
and they can overlap in practice, but each answers a different question about
a program's role.

\begin{center}
\begin{tabular}{p{0.26\textwidth}p{0.60\textwidth}}
\toprule
\textbf{Category} & \textbf{Role in one line} \\
\midrule
System software & Controls the hardware and provides a platform for other software. \\
Application software & Helps the user perform a specific task. \\
Utility software & Performs maintenance and housekeeping tasks. \\
Programming software & Used to write, translate, and test programs. \\
Embedded software & Built into a device to control that device. \\
Middleware & Connects different applications and systems. \\
\bottomrule
\end{tabular}
\end{center}

\section{System software}
\textbf{System software} manages and controls the computer hardware and
provides the \emph{platform} on which application software runs. It works
close to the machine, usually starts when the computer is switched on, and is
not written for one user's task. Common examples are the operating system,
device drivers, firmware such as the BIOS, and utility programs.

The \textbf{operating system (OS)} is the most important system software. It
acts as the interface between the user and the hardware and manages memory,
files, processes, and input/output devices. Windows, Linux, macOS, and Android
are familiar operating systems. Its main functions are process management
(starting, scheduling, and stopping programs), memory management (allocating
and freeing main memory), file management (creating, reading, and organising
files), device management (controlling hardware through drivers), and
providing a user interface. \textbf{Firmware} such as the BIOS or UEFI also
belongs to the system-software layer: it initialises the hardware when the
machine is switched on and then loads the operating system.

A \textbf{device driver} is also system
software: it allows the operating system to communicate with a particular
hardware device. It is easy to mistake a device driver for hardware, but it is
a piece of software (TRAP-10).

\begin{definitionbox}{System software}
Software that controls and coordinates the hardware and provides the services
and platform that other programs need in order to run.
\end{definitionbox}

\section{Application software}
\textbf{Application software} helps the user perform a specific task. It runs
on top of the system software and is written for the user's work rather than
to manage the machine. Word processors, spreadsheets, web browsers, media
players, and accounting packages are all application software.

Application software is often divided into two kinds. \textbf{General-purpose}
software is a ready-made package sold to many users, such as MS Word or MS
Excel; it is also called an \emph{application package}. \textbf{Specific-purpose}
software, also called \emph{tailor-made} or custom software, is built for one
organisation or one job, such as a payroll system or a railway reservation
system.

Application software covers many kinds of user work. Word processors prepare
documents; spreadsheets handle calculations and tables; presentation tools
build slide shows; database packages store and query records; graphics and
multimedia packages edit images, audio, and video; browsers and email clients
handle communication; and accounting packages keep business records. Most
modern application software has a graphical user interface (GUI), so the user
works with windows, icons, and menus instead of typing commands.

\begin{definitionbox}{Application software}
Software that is written to help the user perform a specific task or set of
tasks, running on top of the system software rather than managing the machine.
\end{definitionbox}

\begin{example}
The office desktop runs MS Windows and MS Office. Windows is system software
because it manages the hardware; MS Word and MS Excel are application software
because they help the user create documents and spreadsheets. The same machine
shows all the main categories in use at once.
\end{example}

\section{Utility software}
\textbf{Utility software} is a small program that performs maintenance or
housekeeping tasks on the system. Antivirus programs, disk cleanup tools,
disk defragmentation tools, backup utilities, and file-compression tools are
common examples. Utility software keeps the system healthy, fast, and safe,
and it belongs to the system software family.

A simple way to keep utilities and applications apart is to ask what the
program is for. A utility \emph{maintains the system} (for example, scanning
for viruses or cleaning temporary files), whereas an application helps the
user do \emph{work} (for example, typing a letter or preparing a budget).

\begin{definitionbox}{Utility software}
A small program within the system software family that carries out
maintenance and housekeeping tasks such as virus scanning, disk cleanup,
defragmentation, backup, and compression.
\end{definitionbox}

\begin{example}
A user notices that the office desktop has become slow and that the free disk
space keeps shrinking. Running a disk-cleanup tool to remove temporary files
and a disk-defragmentation tool to reorganise the hard disk restores speed and
space. Both tools are utility software: they maintain the system rather than
help the user produce a document.
\end{example}

\section{Programming software and language translators}
\textbf{Programming software} is used to write, translate, and test computer
programs. The core members are the \textbf{language translators}: a
\textbf{compiler}, an \textbf{interpreter}, and an \textbf{assembler}. A
compiler translates a whole program into machine code at once, while an
interpreter translates and runs the program line by line. An assembler
translates assembly language into machine code.

\begin{center}
\begin{tabular}{p{0.22\textwidth}p{0.64\textwidth}}
\toprule
\textbf{Translator} & \textbf{How it works} \\
\midrule
Compiler & Translates the whole program into machine code before it runs. \\
Interpreter & Translates and executes the program one line at a time. \\
Assembler & Translates assembly language (mnemonic codes) into machine code. \\
\bottomrule
\end{tabular}
\end{center}

Because a compiler translates the whole program first, it reports errors
together and the translated program generally runs faster; because an
interpreter works line by line, it can start running immediately and is
handy for testing small pieces of code.

Around the translators sit other supporting tools: the \textbf{linker}, which
joins separate pieces of a program; the \textbf{loader}, which brings a program
into memory to run; the \textbf{debugger}, which helps find errors; and the
text editor used to write the code. An \textbf{IDE} (Integrated Development
Environment) bundles an editor, a compiler or interpreter, and a debugger into
one package, which makes programming more convenient.

\section{Embedded software}
\textbf{Embedded software} is built into a device in order to control that
device. It is usually stored in read-only memory as \textbf{firmware}, runs
automatically when the device is switched on, and is written for one dedicated
job. Washing machines, microwave ovens, network routers, and the
engine-control units in modern cars all contain embedded software. The user
normally never sees this software directly; it simply makes the device behave
correctly.

\begin{definitionbox}{Embedded software}
Software built into a device and stored as firmware to control the device's
fixed, dedicated functions.
\end{definitionbox}

\section{Middleware}
\textbf{Middleware} is software that connects different applications and
systems so that they can work together. It sits between the operating system
and the application software and provides common services such as messaging,
authentication, and data exchange. An application server that links a web
front end to a back-end database is a typical example. Middleware allows
programs that were built separately to communicate without being rewritten
for each other.

\begin{definitionbox}{Middleware}
Software that sits between the operating system and applications and connects
different programs or systems, providing common services such as messaging,
authentication, and data exchange.
\end{definitionbox}

Middleware is common in large or distributed systems. On a web application,
for example, the browser (client), the web server, and the database are three
different programs. Middleware passing requests and replies between them lets
these separate programs behave as one service, without each program needing to
know the internal details of the others.

\section{Software licensing models}
A software \textbf{licence} states how a program may be used, copied, modified,
and distributed. The main models are summarised below.

\begin{center}
\begin{tabular}{p{0.24\textwidth}p{0.62\textwidth}}
\toprule
\textbf{Model} & \textbf{Key point} \\
\midrule
Proprietary & Source code is kept closed; use is governed by the vendor's licence. \\
Open-source & Source code is available to view, modify, and redistribute under its licence. \\
Free software & Gives the freedom to run, study, modify, and share the software. \\
Freeware & Distributed free of cost, but the source code may not be available. \\
Shareware & Offered as a limited trial; payment is needed for full or continued use. \\
Public domain & Has no copyright and may be used freely for any purpose. \\
Commercial & Sold for profit; usually proprietary. \\
\bottomrule
\end{tabular}
\end{center}

\textbf{Proprietary software} keeps its source code secret and restricts
copying and modification. MS Windows and MS Office are examples.
\textbf{Open-source software} publishes its source code and its licence allows
users to view, change, and share it; Linux, Mozilla Firefox, and LibreOffice
are examples. Note that ``open-source'' refers to \emph{source-code
availability}, not to price. Common open-source licences include the
\textbf{GNU General Public License (GPL)}, the \textbf{Apache License}, and the
\textbf{MIT License}; each sets out the conditions under which the source code
may be used and changed.

Proprietary software is usually supplied under an \textbf{EULA} (End User
License Agreement), which the user accepts before installing the product and
which forbids copying, modifying, or redistributing the program. Because the
source code is closed, only the vendor can improve or fix the software.

\textbf{Free software} is a related but distinct idea: it refers to freedom,
not to price. Free software gives the user the freedom to run the program,
study and modify it, and redistribute copies. It is important to see that both
free software and open-source software make the source code available; they
differ mainly in the philosophy they express, with free software emphasising
the user's freedom. \textbf{Freeware}, by contrast,
simply means free of cost; its source code may remain closed.
\textbf{Shareware} is a trial version distributed free for a limited period,
after which the user must pay for continued or full use. \textbf{Public-domain
software} carries no copyright at all and may be used freely.
\textbf{Commercial software} is sold for profit; it is usually also
proprietary, although a commercial product can occasionally be open-source.

\begin{definitionbox}{Open-source vs freeware vs shareware}
Open-source means the \emph{source code} is available; freeware means the
software is free of \emph{cost} (source may be closed); shareware is a
\emph{trial} that must later be paid for. The three are not the same.
\end{definitionbox}

\begin{example}
An item states, ``Freeware, open-source, and shareware all mean that the
source code is freely available.'' This is wrong. Only open-source guarantees
source-code availability. Freeware is free of cost but may be closed-source,
and shareware is a trial that later requires payment. The wrong option has
merged three different licensing models (TRAP-23).
\end{example}

\section{System software vs application software}
The two broadest groups are worth comparing side by side, because many exam
items turn on exactly this distinction.

\begin{center}
\begin{tabular}{p{0.28\textwidth}p{0.30\textwidth}p{0.30\textwidth}}
\toprule
\textbf{Point} & \textbf{System software} & \textbf{Application software} \\
\midrule
Purpose & Runs and manages the machine & Helps the user do a task \\
Runs & Close to the hardware; starts at boot & On top of the system software \\
Needed for & The computer to work at all & Performing the user's work \\
Examples & Operating system, device driver, BIOS & MS Word, MS Excel, a web browser \\
\bottomrule
\end{tabular}
\end{center}

The same physical computer can run one operating system (system software) and
many applications at the same time. The operating system is what makes the
hardware usable; the applications are what the user actually works with.

\section{Exam retrieval and common confusions}
Five distinctions prevent most errors on this topic.
\begin{enumerate}
  \item \textbf{System vs application}: system software manages the hardware;
    application software helps the user perform a task. The operating system
    is system software, while MS Word is application software.
  \item \textbf{Utility vs application}: a utility \emph{maintains} the system
    (antivirus, disk cleanup); an application helps the user do \emph{work}.
  \item \textbf{A device driver is software}, not hardware; the operating
    system is system software (TRAP-10).
  \item \textbf{Open-source $\neq$ freeware $\neq$ shareware} (TRAP-23):
    open-source gives the source code, freeware is free of cost, and shareware
    is a paid trial.
  \item \textbf{Free software is about freedom, not price}; it should not be
    confused with freeware, which is about price alone.
\end{enumerate}

\begin{example}
An item asks which of a compiler, an antivirus program, MS Excel, and the
operating system is \emph{not} system software. Trace the categories: a
compiler is programming software (system family), an antivirus program is
utility software (system family), and the operating system is system software.
MS Excel is application software, so it is the odd one out.
\end{example}

\subsection*{Exercises}
\begin{enumerate}[label=\arabic*.]
  \item Define software and distinguish it from hardware.
  \item List the six main categories of software and give one example of each.
  \item Distinguish system software from application software, with one
    example of each.
  \item Explain the difference between utility software and application
    software.
  \item Distinguish open-source software, freeware, and shareware, and state
    which one guarantees access to the source code.
  \item Classify each of the following as system, application, utility,
    programming, embedded, or middleware software: MS Word; an antivirus
    program; a compiler; the firmware in a router; an application server.
\end{enumerate}

\subsection*{Solutions}
\begingroup\small
\begin{enumerate}[label=\arabic*.]
  \item Software is a set of programs (instructions) that tells the computer
    what to do; hardware is the physical equipment that runs those
    instructions.
  \item System software (operating system), application software (MS Word),
    utility software (antivirus), programming software (compiler), embedded
    software (washing-machine firmware), and middleware (application server).
  \item System software manages and controls the hardware and provides a
    platform; application software helps the user perform a specific task.
    The operating system is system software; MS Excel is application
    software.
  \item Utility software performs maintenance and housekeeping tasks such as
    virus scanning or disk cleanup; application software helps the user carry
    out work such as writing a document or preparing a budget.
  \item Open-source software publishes its source code and allows it to be
    viewed, modified, and shared; freeware is free of cost but may keep its
    source code closed; shareware is a trial that later requires payment.
    Only open-source software guarantees access to the source code.
  \item MS Word is application software; an antivirus program is utility
    software; a compiler is programming software; the firmware in a router is
    embedded software; and an application server is middleware.
\end{enumerate}
\endgroup
\end{document}
